% Generated draft: requires mathematical and source review.
\documentclass[11pt]{article}
\usepackage[margin=25mm]{geometry}
\usepackage{fontspec}
\setmainfont{Latin Modern Roman}
\usepackage{amsmath,amssymb,mathtools,unicode-math}
\setmathfont{Latin Modern Math}
\usepackage{enumitem,needspace,xurl,hyperref}
\setlength{\parindent}{0pt}
\setlength{\parskip}{5pt}
\setlength{\emergencystretch}{3em}
\newcommand{\N}{\mathbb{N}}
\newcommand{\Z}{\mathbb{Z}}
\newcommand{\Q}{\mathbb{Q}}
\newcommand{\R}{\mathbb{R}}
\newcommand{\C}{\mathbb{C}}
\newcommand{\F}{\mathbb{F}}
\newcommand{\A}{\mathbb{A}}
\newcommand{\PP}{\mathbb{P}}
\newcommand{\E}{\mathbb{E}}
\newcommand{\eps}{\varepsilon}
\newcommand{\dd}{\,\mathrm{d}}
\begin{document}
\section*{UnsolvedMath ID: 30000027}
\textbf{Source problem code:} catalogue entry
\subsection{Definitions and mathematical statement}
Consider nonlinear wave equations in three spatial dimensions, with unknown $u=u(t,x)$ and spatial variable $x\in\Omega\subset\R^3$ in the domain intended by the source. The indicated nonlinearity depends on the jet
\[
 f\bigl(u_t,\nabla_xu,\nabla_xu_t,\nabla_x^2u\bigr),
\]
where $\nabla_x^2u=(\partial_{x_i}\partial_{x_j}u)_{i,j=1}^3$. ``Behavior near zero'' concerns these arguments tending jointly to zero, and may involve their vanishing orders or other local asymptotics; no specific order is prescribed.

The question is how effective versus ineffective damping changes the available $L^p$-to-$L^q$ decay estimates and how this depends on that local behavior of $f$, including the variable-coefficient and nonlinear settings. Here an $L^p$-to-$L^q$ estimate controls a solution quantity in spatial $L^q$ at time $t$ by appropriate initial-data norms based on $L^p$, with a rate as $t\to\infty$. The target is to determine the admissible exponents, decay rates and assumptions in the damping regimes, not a supplied numerical rate or universal inequality.

The extract does not state a particular evolution operator, damping coefficient, spatial domain or boundary condition, initial-data regularity and size, or whether the decaying quantity is $u$, its derivatives, or energy. It also does not define effective and ineffective damping for the intended equation class. Those regimes depend on the damping law and cannot be specified universally by importing a threshold from one model. A fully quantified decay theorem therefore requires this missing model data; no constant-coefficient equation or small-data hypothesis is silently substituted.

The literature assessment checked on 21 August 2026 reports substantial decay and asymptotic theory for important effective and noneffective damping classes, but says no general three-dimensional relation covering the variable-coefficient and nonlinear scope in this question was found. This is a dated partial-progress assessment requiring independent verification.
\subsection{Short English statement}
Determine how damping regimes and the near-zero structure of the displayed nonlinearity govern three-dimensional wave decay from $L^p$ data to $L^q$ solution norms. The equation class and precise damping thresholds must be specified before a universal rate can be formulated.
\subsection{Sources}
\begin{itemize}
\item \emph{Wave Motion} (2004), Oberwolfach Reports, DOI 10.4171/owr/2004/05, pp. 243--294, \url{https://doi.org/10.4171/owr/2004/05}.
\item Literature assessment checked 21 August 2026: ``Oberwolfach report,'' \url{https://ems.press/content/serial-article-files/45893}; ``Recent evolution-equation paper,'' \url{https://arxiv.org/abs/2412.10022}.
\item Local UnsolvedMath registry, permanent ID $30000027$, exact source code \texttt{catalogue entry}, created 25 August 2026 and updated 16 September 2026. Its clean statement is labelled corrected and verified against the report; partially-solved is the supplied dated literature label.
\end{itemize}
\end{document}
